%========================================================

\begin{frame}{Dispersione cromatica}
    \begin{itemize}
        \item L'indice di rifrazione $n$ non è una costante assoluta, ma dipende dalla lunghezza d'onda $\lambda$ della luce che attraversa il mezzo ($n = n(\lambda)$)
        \item Poichè $v=\frac{c}{n}$, le diverse frequenze (colori) viaggiono a velocità differenti all'interno del prisma.
        \item Per la Legge di Snell ($n_1sin\theta_1 = n_2sin\theta_2$), angoli di rifrazione diversi portano alla scomposizione della luce bianca nello spettro visibile.
        \item In genere, per i materiali trasparenti, l'indice $n$ aumenta al diminuire di $\lambda$ (il violetto viene deviato più del rosso)
    \end{itemize}
\end{frame}

\begin{frame}{Diffrazione}
    \begin{itemize}
        \item \textbf{Principio di Huygens}: ogni punto di un fronte d'onda si comporta come una sorgente di onde sferiche secondarie
        \item Il fenomeno è rilevante quando la dimensione dell'apertura $a$ è paragonabile alla lunghezza d'onda $\lambda$
        \item (Fenditura singola):$$a \sin \theta = m\lambda \quad (m = \pm 1, \pm 2, \dots)$$
        \item Le ombre non sono mai perfettamente nette; la luce invade la zona d'ombra geometrica
    \end{itemize}
\end{frame}


\begin{frame}{Interferenza (L'Esperimento di Young)}
    \begin{itemize}
        \item Due sorgenti coerenti separate da una distanza $d$ producono frange luminose e scure sullo schermo
        \item Differenza di cammino ottico: $\Delta r = d \sin \theta$.
        \item Condizioni matematiche:
        \begin{itemize}
            \item Costruttiva (Luce): $\Delta r = m\lambda$
            \item Distruttiva (Buio): $\Delta r = (m + \frac{1}{2})\lambda$
        \end{itemize}
        \item \textbf{Interfrangia}: La distanza tra due massimi consecutivi è $\Delta y = \frac{\lambda L}{d}$
    \end{itemize}
\end{frame}

\begin{frame}{Polarizzazione}
    \begin{itemize}
        \item \textbf{Natura trasversale:} Il campo elettrico $\vec{E}$ della luce oscilla perpendicolarmente alla direzione di propagazione
        \item \textbf{Filtro Polarizzatore:} Seleziona un'unica direzione di oscillazione, assorbendo le altre
        \item \textbf{Legge di Malus:}  Descrive l'intensità della luce $I$ che attraversa un secondo polarizzatore (analizzatore) inclinato di un angolo $\theta$: $$I = I_0 \cos^2 \theta$$
        \item \textbf{Applicazioni:} Occhiali da sole polarizzati, schermi LCD e fotografia
    \end{itemize}
\end{frame}